\documentclass[11pt]{article}
\usepackage[margin=1in]{geometry}
\usepackage{amsmath,amssymb,amsthm,mathtools}
\usepackage{graphicx,microtype}
\usepackage[hidelinks]{hyperref}

\newtheorem{theorem}{Theorem}
\newtheorem{lemma}[theorem]{Lemma}
\newtheorem{remark}[theorem]{Remark}
\newcommand{\R}{\mathbb R}
\newcommand{\cE}{\mathcal E}
\newcommand{\cD}{\mathcal D}

\title{Quantum Confinement at Generic Tangency Points}
\author{Autonomous mathematical discovery run \texttt{fa\_banach\_001}}
\date{13 August 2026}

\begin{document}
\maketitle

\begin{abstract}
For a compact generic two-dimensional almost-Riemannian structure, we prove
that the ordinary Laplace--Beltrami operator on every connected component of
the regular set is essentially self-adjoint, including in the presence of
isolated tangency points.  This answers the tangency-point question recorded
in arXiv:2011.03300.  The key is a pair of scale-sharp Hardy estimates for the
quadratic tangency normal form: a transverse constant-three inequality on
the cusp side and a constant-$9/4$ weighted inequality on the opposite side.
Together with a complete logarithmic sector and the known Grushin estimate,
these give the Agmon barrier after an IMS partition.
\end{abstract}

\section{Question and result}

Beschastnyi, Boscain, and Pozzoli ask on page 5 of
\cite{BBP} whether the ordinary Laplace--Beltrami operator remains
essentially self-adjoint when the two-step hypothesis is removed and
tangency points occur:
\begin{quote}
``In presence of tangency points even the essential self-adjointness of the
standard Laplace--Beltrami operator (without the term $-cK$) is an open
question.''
\end{quote}
The source passage is reproduced in Figure~\ref{fig:question}.

\begin{figure}[ht]
  \centering
  \includegraphics[width=.94\linewidth]{figures/open_problem_crop.png}
  \caption{The open question on page 5 of arXiv:2011.03300v2.}
  \label{fig:question}
\end{figure}

A generic $2$-ARS has a smooth singular curve, ordinary Grushin points, and
finitely many isolated quadratic tangency points.  At a tangency, after a
smooth change of coordinates and an orthonormal change of frame, its local
normal form is
\begin{equation}\label{eq:normal}
 X_1=\partial_x,\qquad
 X_2=f(x,y)\partial_y,\qquad
 f=(y-x^2\psi(x))e^{\Psi(x,y)},
\end{equation}
where $\psi(0)>0$ and $\Psi$ is smooth; see \cite{ABCGS,Beschastnyi}.
The regular-set metric and volume are
\[
 g=dx^2+f^{-2}dy^2,\qquad d\omega=|f|^{-1}dx\,dy,
\]
and the nonnegative minimal operator $H=-\Delta_g$ has form
\begin{equation}\label{eq:form}
 \cE[u]=\int\bigl(|\partial_xu|^2+f^2|\partial_yu|^2\bigr)\,d\omega,
 \qquad u\in C_c^\infty(\Omega).
\end{equation}

\begin{theorem}[Generic tangencies preserve confinement]\label{thm:main}
Let $M$ be a compact generic two-dimensional almost-Riemannian manifold,
and let $Z$ be its singular set.  For every connected component
$\Omega$ of $M\setminus Z$, the Laplace--Beltrami operator
\[
 -\Delta_g:C_c^\infty(\Omega)\longrightarrow L^2(\Omega,\omega)
\]
is essentially self-adjoint.  In particular, isolated quadratic tangency
points require no boundary conditions and do not permit quantum transport
between components of $M\setminus Z$.
\end{theorem}

\section{Proof intuition}

Away from tangencies this is the established Grushin theorem
\cite{BL,FPR}.  Near a tangency we may discard smooth lower-order
distortions and inspect $f=y-x^2$.  Above the parabola, every horizontal
slice is the shrinking interval $(-\sqrt y,\sqrt y)$.  A one-dimensional
completion of squares on the scaled interval produces a Hardy barrier with
constant $3$.  It is stronger than the squared gradient of the natural cusp
gauge in the central cone; the two edge cones are ordinary Grushin collars.

Below the parabola there are two different ends.  The sector $-y\gtrsim
x^2$ is complete in the logarithmic variable $-\log(-y)$.  In the remaining
horizontal sectors, the density is asymptotic to $|x|^{-2}$, and another
completion of squares gives the constant $9/4>1$ required by the Agmon
argument.  A partition depending slowly on the logarithmic angular ratio
$(-y)/x^2$ has arbitrarily small IMS cost.  Both strict inequalities survive
the smooth factors in \eqref{eq:normal}.

\section{Two exact model identities}

We first record the elementary algebra responsible for the strict barriers.

\begin{lemma}[Upper cusp]\label{lem:upper}
Let $q=y-x^2$ and $\Omega_+=\{q>0\}$.  For every
$u\in C_c^\infty(\Omega_+)$,
\begin{equation}\label{eq:upperhardy}
 \int_{\Omega_+}\frac{|u_x|^2}{q}\,dx\,dy
 \ge 3\int_{\Omega_+}\frac{y|u|^2}{q^3}\,dx\,dy.
\end{equation}
If $\rho_+=q/\sqrt y$, then the right side is
$3\int |u|^2\rho_+^{-2}\,d\omega$ and, with $z=x^2/y$,
\begin{equation}\label{eq:uppergrad}
 |\nabla\rho_+|_g^2
 =\frac14\{y(1-z^2)^2+16z\}.
\end{equation}
Consequently, on every central cone $z\le z_0<3/4$ and in a sufficiently
small neighborhood of the tangency,
$|\nabla\rho_+|^2\le3-\varepsilon$ for some $\varepsilon>0$.
\end{lemma}

\begin{proof}
For fixed $y>0$, put $x=\sqrt y\,r$ and $U(r)=u(\sqrt y r,y)$.
For $w=(1-r^2)^{-1}$ and $a=3r/(1-r^2)$, integration of the square gives
\begin{align*}
 \int_{-1}^1w|U'|^2,dr
 &=\int_{-1}^1w|U'+aU|^2,dr
   +\int_{-1}^1\bigl((wa)'-wa^2\bigr)|U|^2,dr\\
 &=\int_{-1}^1w|U'+aU|^2,dr
   +3\int_{-1}^1\frac{|U|^2}{(1-r^2)^3},dr.
\end{align*}
Scaling and integrating in $y$ proves \eqref{eq:upperhardy}.  Direct
differentiation gives \eqref{eq:uppergrad}.
\end{proof}

\begin{lemma}[Lower horizontal cones]\label{lem:lower}
Let $b=-y\ge0$, $q=x^2+b$, and work on one half-plane $x>0$ or $x<0$.
For compactly supported $u$ and every fixed $b$,
\begin{equation}\label{eq:lowerhardy}
 \int\frac{|u_x|^2}{x^2+b},dx
 \ge\int \frac{3(3x^2-b)}{4(x^2+b)}
              \frac{|u|^2}{x^2}\frac{dx}{x^2+b}.
\end{equation}
Thus on $0\le b\le x^2/4$ the coefficient multiplying
$x^{-2}d\omega$ is at least $33/20>1$.
\end{lemma}

\begin{proof}
Use $w=(x^2+b)^{-1}$ and $a=-3/(2x)$.  The same square identity gives
\[
 (wa)'-wa^2
 =\frac{3(3x^2-b)}{4x^2(x^2+b)^2},
\]
which is \eqref{eq:lowerhardy}.  The stated minimum is attained at
$b=x^2/4$.
\end{proof}

\section{Agmon localization at the tangency}

We use the following standard boundary-local form of the Agmon argument.
It is included to make clear exactly how the two strict constants enter.

\begin{lemma}[Hybrid end lemma]\label{lem:agmon}
Let $H$ be the nonnegative divergence-form operator associated with a
strongly local form $\cE$ on the interior of a manifold.  Suppose a boundary
end has a finite smooth IMS partition $\sum_j\chi_j^2=1$ with the following
properties.
\begin{enumerate}
\item Each localized end is either complete (it has compactly supported
cutoffs converging to one with gradient converging uniformly to zero), or it
has a positive boundary gauge $\rho_j$ and constants $\varepsilon,C>0$ such
that
\[
 \cE[v]\ge\int
 (|\nabla\rho_j|^2+\varepsilon)\rho_j^{-2}|v|^2,d\omega-C\|v\|_2^2
\]
for functions supported in that localized end.
\item The IMS error $\sum_j|\nabla\chi_j|^2$ is absorbed by the strict
$\varepsilon\rho_j^{-2}$ terms on transition regions.
\end{enumerate}
If all remaining boundary patches are already limit point, then $H$ is
essentially self-adjoint at this end.
\end{lemma}

\begin{proof}
For completeness, let $h\in\ker(H^*+\lambda)$ with $\lambda>C+1$ after a
harmless spectral shift.  For a compactly supported Lipschitz cutoff $F$,
the weak equation and the product rule give the Agmon identity
\begin{equation}\label{eq:agmon}
 \cE[Fh]+\lambda\|Fh\|_2^2
 =\int |\nabla F|^2|h|^2,d\omega.
\end{equation}
On complete pieces take the usual slowly varying exhaustion.  On a gauged
piece take the standard power cutoff in $\rho_j$, with exponent chosen
strictly between $1$ and
$\sqrt{1+\varepsilon/\|\nabla\rho_j\|_infty^2}$, and then let its inner
radius tend to zero.  The gauge inequality absorbs the right side of
\eqref{eq:agmon}; the remaining innermost annulus tends to zero by
$h\in L^2$.  Applying \eqref{eq:agmon} to $F\chi_jh$, summing, and using the
IMS identity leaves precisely the transition error in item 2, which is also
absorbed.  Hence every localized $h$ vanishes, so $h=0$.  The usual
deficiency criterion for a semibounded symmetric operator proves essential
self-adjointness.  This is the localized version of the weak-Hardy/Agmon
proof in \cite{PRS,FPR}.
\end{proof}

\begin{proof}[Proof of Theorem~\ref{thm:main}]
The assertion is already known near every non-tangency point of $Z$ by
\cite{BL,FPR}.  It remains to verify Lemma~\ref{lem:agmon} near one isolated
tangency.  Normalize \eqref{eq:normal} so that $\psi(0)=1$.  We first use the
canonical model $\Psi=0$, $\psi=1$.

On $\Omega_+=\{y>x^2\}$ choose a central parabolic cutoff supported where
$x^2/y<z_0<3/4$.  Lemma~\ref{lem:upper} gives the gauged estimate with a
strict margin.  The two complementary edge collars are ordinary Grushin
ends; equivalently, $q/(2|x|)$ is their regular normal variable, and the
critical Grushin estimate of \cite{BL,FPR} is uniform there.  Choose the
angular partition as a smooth function of $r=x/\sqrt y$.  Its IMS cost is
$O(y^{-1})$.  On every transition strip $|r|\le r_1<1$,
$\rho_+^{-2}=y^{-1}(1-r^2)^{-2}$, so the strict part of
\eqref{eq:upperhardy} absorbs this cost (one may spread the angular
transition before fixing the neighborhood).  This proves the upper-side
condition.

For the lower component, the part with $y>0$ consists of two ordinary
Grushin collars.  Put $b=-y$ on $y\le0$.  In the vertical cone
$b\ge\kappa x^2$, the function $s=-\log b$ tends to infinity at the tip and
\[
 |\nabla s|^2=(x^2+b)^2/b^2\le(1+\kappa^{-1})^2.
\]
Thus this is a complete end after using cutoffs whose transition length in
$s$ tends to infinity.  In the horizontal cones $b\le x^2/4$,
Lemma~\ref{lem:lower} gives a gauge $\rho_-=|x|$ with
$|\nabla\rho_-|=1$ and Hardy coefficient at least $33/20>1$.
Between the cones take a partition depending on
$\log(b/x^2)$.  By choosing the two angular thresholds with a large
logarithmic separation, its $x^{-2}$ part is arbitrarily small; its bounded
part is harmless after the spectral shift.  Hence its IMS cost is absorbed
by the $13/20$ margin in Lemma~\ref{lem:lower}.  The lower side also satisfies
Lemma~\ref{lem:agmon}.

Finally restore $\psi$ and $\Psi$.  Shrinking the coordinate neighborhood,
$\psi(x)=1+O(x)$ and $e^{\Psi}=e^{\Psi(0,0)}(1+O(|x|+|y|))$.  A constant
linear rescaling normalizes $e^{\Psi(0,0)}$.  The form, both gauges, and all
square remainders are therefore perturbed by relative $o(1)$ terms plus
bounded lower-order terms.  The upper central margin, the lower coefficient
$33/20>1$, and the logarithmically adjustable IMS margins persist.  This
also proves uniformity over the finitely many tangencies.  Lemma
\ref{lem:agmon}, followed by a finite global IMS partition on compact $M$,
completes the proof.
\end{proof}

\section{Scope, verification, and novelty}

The theorem addresses the standard generic case: tangencies are isolated and
have the quadratic normal form \eqref{eq:normal}.  It does not cover
arbitrarily flat non-generic contacts or singular sets that cease to be
embedded curves.

The exact algebra in Lemmas~\ref{lem:upper} and \ref{lem:lower}, as well as
formula \eqref{eq:uppergrad}, was independently checked by the SymPy script
\texttt{code/verify\_tangency\_hardy.py}.  The script is not a proof of the
localization lemma.

Focused searches through 13 August 2026 used the phrases ``tangency points
essential self-adjointness Laplace--Beltrami'', ``generic 2D
almost-Riemannian quantum confinement'', and the canonical model
$X_2=(y-x^2)\partial_y$.  The later paper \cite{Beschastnyi} develops a Lie
groupoid method that handles inverse-square perturbations at tangencies, but
explicitly says polynomial weights do not settle the bare Laplacian there.
The $\alpha$-Grushin calculus of \cite{BQ} treats non-tangent singular
hypersurfaces.  No statement of Theorem~\ref{thm:main} was found.  This is a
bounded novelty search, not an exhaustive bibliographic claim.

The recommended human check is the uniform IMS/Agmon passage in
Lemma~\ref{lem:agmon} and the final smooth-normal-form stability paragraph.
The two decisive constants are exact and have ample strict margins in the
regions where they are used.

\begin{thebibliography}{9}
\bibitem{BBP}
I. Beschastnyi, U. Boscain, and E. Pozzoli,
\emph{Quantum confinement for the curvature Laplacian $-\Delta+cK$ on
2D-almost-Riemannian manifolds}, arXiv:2011.03300v2 (2021).

\bibitem{BL}
U. Boscain and C. Laurent,
\emph{The Laplace--Beltrami operator in almost-Riemannian geometry},
Ann. Inst. Fourier 63 (2013), 1739--1770; arXiv:1105.4687.

\bibitem{FPR}
V. Franceschi, D. Prandi, and L. Rizzi,
\emph{On the essential self-adjointness of singular sub-Laplacians},
Potential Anal. 53 (2020), 89--112; arXiv:1708.09626.

\bibitem{PRS}
D. Prandi, L. Rizzi, and M. Seri,
\emph{Quantum confinement on non-complete Riemannian manifolds},
J. Spectr. Theory 8 (2018), 1221--1280; arXiv:1609.01724.

\bibitem{ABCGS}
A. Agrachev, U. Boscain, G. Charlot, R. Ghezzi, and M. Sigalotti,
\emph{Two-dimensional almost-Riemannian structures with tangency points},
Ann. Inst. H. Poincar\'e Anal. Non Lin\'eaire 27 (2010), 793--807.

\bibitem{Beschastnyi}
I. Beschastnyi,
\emph{Closure of the Laplace--Beltrami operator on 2D almost-Riemannian
manifolds and semi-Fredholm properties of differential operators on Lie
manifolds}, arXiv:2104.07745 (2021).

\bibitem{BQ}
I. Beschastnyi and H. Quan,
\emph{Self-adjointness criteria and self-adjoint extensions of the
Laplace--Beltrami operator on $\alpha$-Grushin manifolds},
arXiv:2305.08280 (2023).
\end{thebibliography}

\end{document}
